% Options for packages loaded elsewhere
\PassOptionsToPackage{unicode}{hyperref}
\PassOptionsToPackage{hyphens}{url}
\documentclass[
]{article}
\usepackage{xcolor}
\usepackage{amsmath,amssymb}
\setcounter{secnumdepth}{-\maxdimen} % remove section numbering
\usepackage{iftex}
\ifPDFTeX
  \usepackage[T1]{fontenc}
  \usepackage[utf8]{inputenc}
  \usepackage{textcomp} % provide euro and other symbols
\else % if luatex or xetex
  \usepackage{unicode-math} % this also loads fontspec
  \defaultfontfeatures{Scale=MatchLowercase}
  \defaultfontfeatures[\rmfamily]{Ligatures=TeX,Scale=1}
\fi
\usepackage{lmodern}
\ifPDFTeX\else
  % xetex/luatex font selection
\fi
% Use upquote if available, for straight quotes in verbatim environments
\IfFileExists{upquote.sty}{\usepackage{upquote}}{}
\IfFileExists{microtype.sty}{% use microtype if available
  \usepackage[]{microtype}
  \UseMicrotypeSet[protrusion]{basicmath} % disable protrusion for tt fonts
}{}
\makeatletter
\@ifundefined{KOMAClassName}{% if non-KOMA class
  \IfFileExists{parskip.sty}{%
    \usepackage{parskip}
  }{% else
    \setlength{\parindent}{0pt}
    \setlength{\parskip}{6pt plus 2pt minus 1pt}}
}{% if KOMA class
  \KOMAoptions{parskip=half}}
\makeatother
\usepackage{longtable,booktabs,array}
\usepackage{caption}
\captionsetup[table]{skip=6pt}
\newcounter{none} % for unnumbered tables
\usepackage{calc} % for calculating minipage widths
% Correct order of tables after \paragraph or \subparagraph
\usepackage{etoolbox}
\makeatletter
\patchcmd\longtable{\par}{\if@noskipsec\mbox{}\fi\par}{}{}
\makeatother
\setlength{\emergencystretch}{3em} % prevent overfull lines
\providecommand{\tightlist}{%
  \setlength{\itemsep}{0pt}\setlength{\parskip}{0pt}}
\usepackage{bookmark}
\IfFileExists{xurl.sty}{\usepackage{xurl}}{} % add URL line breaks if available
\urlstyle{same}
% fallback for those not using the hyperref driver hyperxmp:
\makeatletter
\@ifundefined{xmpquote}{\newcommand{\xmpquote}[1]{#1}}{}
\makeatother
\hypersetup{
  hidelinks,
  pdfcreator={LaTeX via pandoc}}

\author{}
\date{}

\begin{document}

\section{Knowledge Distillation and Deployment of a Lightweight
Command-Parsing LM on a Local
Cloud}\label{knowledge-distillation-and-deployment-of-a-lightweight-command-parsing-lm-on-a-local-cloud}

\textbf{Name:} Ritika\\
\textbf{Roll number:} 230101067\\
\textbf{Course:} CS3081 Cloud Computing\\
\textbf{Repository:}
\texttt{https://github.com/\textless{}your-username\textgreater{}/230101067\_CloudComputing\_Assignment}

\subsection{1. Architecture and design
decisions}\label{architecture-and-design-decisions}

The project is a local cloud with four layers. A k3d cluster is created
from \texttt{cluster/k3d.yaml} (one server, one agent, Traefik and
metrics-server kept enabled). SeaweedFS runs in namespace \texttt{data}
and exposes an S3-compatible API. A parallel Kubernetes Indexed Job in
namespace \texttt{pipeline} reads
\texttt{datasets/commands\_pool.jsonl}, labels each command with a
host-run Ollama \texttt{phi3.5}, and writes one raw shard object per
worker. A cleaning step validates the replies and writes
\texttt{datasets/v1/train.jsonl} with a SHA-256 checksum. Training
distils the teacher replies into \texttt{EleutherAI/pythia-70m} and
stores \texttt{model.tar.gz} and \texttt{model.sha256} in the
\texttt{models} bucket. The FastAPI service runs in namespace
\texttt{app}, downloads the model at start-up, and is exposed by a
Traefik Ingress.

The main design choice is to keep weights out of the image and out of
Git. The service image contains only Python and pinned libraries;
\texttt{/readyz} stays 503 until the model has been downloaded,
checksum-verified, extracted, and loaded. This makes cold start
observable and makes a model update a change of
\texttt{MODEL\_URI}/\texttt{MODEL\_VERSION}, not a rebuild. Ollama runs
on the host rather than in the cluster because the teacher is only
needed during Task 1 and host networking avoids a multi-gigabyte GPU
image. The workers reach it through \texttt{host.k3d.internal:11434}.
Every component is stateless except SeaweedFS, so Pod deletion and
rolling updates are safe.

Version pins live in \texttt{training/config.yaml},
\texttt{service/requirements.txt}, \texttt{worker/requirements.txt} and
\texttt{training/requirements.txt}. All cluster configuration is in
\texttt{k8s/}; the only host-generated state is the git-ignored
\texttt{.env.local} containing S3 credentials.

\subsection{2. Task 1: object storage and parallel teacher
labelling}\label{task-1-object-storage-and-parallel-teacher-labelling}

\subsubsection{2.1 Storage layout}\label{storage-layout}

SeaweedFS is deployed by \texttt{k8s/seaweedfs.yaml} into namespace
\texttt{data}. Two buckets are created: \texttt{datasets} and
\texttt{models}. \texttt{scripts/init\_storage.py} uploads the provided
\texttt{commands\_pool.jsonl} to \texttt{datasets/commands\_pool.jsonl}.
\texttt{prompts/teacher\_prompt.txt} asks for exactly the four keys and
the eight allowed intents, gives three examples, and is fixed at
\texttt{temperature=0}, \texttt{seed=42}, \texttt{top\_k=1}.

\subsubsection{2.2 Sharding and
resumability}\label{sharding-and-resumability}

\texttt{worker/worker.py} reads \texttt{SHARD\_INDEX} and
\texttt{SHARD\_COUNT}, keeps the command on line \texttt{i} when
\texttt{i\ \%\ SHARD\_COUNT\ ==\ SHARD\_INDEX}, and writes
\texttt{datasets/teacher/shard-\textless{}index\textgreater{}.jsonl}
with exactly \texttt{\{"id","text","teacher\_raw"\}}. The Kubernetes Job
uses \texttt{completionMode:\ Indexed}, \texttt{completions:\ 4},
\texttt{parallelism:\ 2}; Kubernetes supplies
\texttt{JOB\_COMPLETION\_INDEX}, which the worker exports as
\texttt{SHARD\_INDEX}. With four completion indexes, each of the 3,000
line numbers is selected by exactly one worker, so there are no gaps or
repeats.

Resume is built in: before labelling, the worker reads its own existing
shard object into a set of completed ids and skips them. Progress is
written after every reply by default (atomic temp-file + rename), so an
interrupted multi-hour run loses at most one command. The cleaner also
refuses to proceed if ids are repeated across shards;
\texttt{-\/-allow-incomplete} is only for smoke runs.

\subsubsection{2.3 Cleaning rules}\label{cleaning-rules}

\texttt{worker/clean.py} applies rejection precedence (a)
-\textgreater{} (b) -\textgreater{} (c) -\textgreater{} (d) and counts
each record exactly once:

\begin{itemize}
\tightlist
\item
  \begin{enumerate}
  \def\labelenumi{(\alph{enumi})}
  \tightlist
  \item
    \texttt{teacher\_raw} is not valid JSON or is not an object;
  \end{enumerate}
\item
  (b) keys are not exactly \texttt{intent}, \texttt{date},
  \texttt{time}, \texttt{place};
\item
  (c) \texttt{intent} is not one of the eight allowed values;
\item
  (d) a slot is neither null nor a string that appears in the command
  (case-insensitive, whitespace-normalised);
\item
  duplicates are removed by normalised command text, keeping the
  smallest id.
\end{itemize}

Survivors are sorted by numeric id and written to
\texttt{datasets/v1/train.jsonl}; the SHA-256 line is written to
\texttt{datasets/v1/train.sha256}. The accounting invariant below is
asserted in code:

\begin{verbatim}
input = invalid_json + wrong_keys + bad_intent + bad_slot + duplicates + kept
\end{verbatim}

\texttt{\textless{}FILL:\ Table\ 1\ -\ Task\ 1\ counts\ for\ the\ smoke\ run\ and\ the\ full\ run\ if\ available\textgreater{}}

\subsubsection{2.4 Sharding rationale}\label{sharding-rationale}

Four shards are the smallest count that gives two waves at
\texttt{parallelism:\ 2}, so the cluster exercises both parallel workers
and resume without creating idle pods. Sharding by line number, rather
than by id, is exactly what the assignment specifies and is independent
of the order of ids. Because every worker owns a disjoint residue class,
the union is complete without a merge step. The cost is that a worker's
commands are interleaved through the pool, but the worker sorts its own
output by id, so each object is deterministic. Four workers also mean
four objects, which is small enough to inspect by hand and large enough
to show that the cleaning step really reads multiple shards. The local
smoke run used the same code with four shards and one worker at a time,
for 50 commands per shard (200 total), because a full CPU teacher run of
3,000 commands takes several hours; the Kubernetes Job is the intended
full path and is resumable.

\subsection{3. Task 2: training the
student}\label{task-2-training-the-student}

\subsubsection{3.1 Configuration}\label{configuration}

\texttt{training/config.yaml} records the six required items:

{\def\LTcaptype{none} % do not increment counter
\begin{longtable}[]{@{}
  >{\raggedright\arraybackslash}p{(\linewidth - 2\tabcolsep) * \real{0.5000}}
  >{\raggedright\arraybackslash}p{(\linewidth - 2\tabcolsep) * \real{0.5000}}@{}}
\toprule\noalign{}
\begin{minipage}[b]{\linewidth}\raggedright
Item
\end{minipage} & \begin{minipage}[b]{\linewidth}\raggedright
Value
\end{minipage} \\
\midrule\noalign{}
\endhead
\bottomrule\noalign{}
\endlastfoot
Random seed & 42 (\texttt{torch.manual\_seed}, NumPy, Python
\texttt{random}; fixed in code) \\
Learning rate & 5e-5 \\
Batch size & 4 (gradient accumulation 4, effective 16) \\
Epochs & 2 \\
Maximum sequence length & 256 \\
Library versions & Python 3.12, torch 2.4.1, transformers 4.44.2,
tokenizers 0.19.1, safetensors 0.4.5, huggingface-hub 0.24.6, numpy
1.26.4, PyYAML 6.0.2, boto3 1.35.24 \\
\end{longtable}
}

\texttt{training/config\_v2.yaml} changes only the seed (43), epochs
(3), version and loss file. \texttt{make\ train} reads only
\texttt{training/config.yaml}; \texttt{make\ train\ SUBSET=50} reads the
same file and caps the rows. \texttt{training/loss.csv} is written and
flushed after every epoch.

\subsubsection{3.2 Example format and
loss}\label{example-format-and-loss}

Each record becomes
\texttt{Command:\ \textless{}text\textgreater{}\textbackslash{}nJSON:\ \textless{}teacher\ reply\textgreater{}}.
The prompt tokens are masked with \texttt{-100}, so the cross-entropy is
computed only on the teacher reply tokens plus EOS; this matches ``the
model learns to produce the teacher's reply, given the command'' and
prevents the 70M model from spending capacity on reproducing the prompt.
The model is \texttt{EleutherAI/pythia-70m}, loaded in float32, trained
with AdamW and a linear warm-up schedule on CPU. The target is
sequence-level: the exact teacher string, not logits.

\texttt{\textless{}FILL:\ Table\ 2\ -\ v1\ and\ v2\ loss\ per\ epoch,\ records\ used,\ wall\ time\textgreater{}}

\subsubsection{3.3 Evaluation}\label{evaluation}

Predictions are generated for the teacher (\texttt{phi3.5}, same prompt,
temperature 0, seed 42), untrained Pythia-70M, and the two students. Raw
outputs and score reports are in \texttt{results/}. The official
\texttt{tools/score.py} computes JSON validity, intent accuracy, slot
accuracies and full-frame accuracy.

\texttt{\textless{}FILL:\ Table\ 3\ -\ teacher/base/v1/v2\ scores\textgreater{}}

\texttt{\textless{}FILL:\ short\ discussion\ of\ what\ changed\textgreater{}}

\subsubsection{3.4 Sequence-level versus logit-based
distillation}\label{sequence-level-versus-logit-based-distillation}

Sequence-level distillation trains the student on the teacher's
generated sequence, treating it as a hard target. The loss is the
ordinary next-token cross-entropy on that sequence. Logit-based
distillation trains the student to match the teacher's full next-token
probability distribution at each position, typically with a
temperature-softened KL divergence. Logit-based distillation is not
possible here because the teacher is served by Ollama, which exposes
only the generated text through \texttt{/api/generate}; it does not
return per-token logits. Even if it did, the teacher and student use
different tokenizers, so their vocabularies would need an alignment
step. A limitation of the sequence-level approach is that it discards
the teacher's uncertainty and cannot teach the student about near-miss
alternatives; with only about 1,500 examples, the student can also
memorise surface forms instead of generalising to unseen commands.

\subsection{4. Task 3: inference service and
container}\label{task-3-inference-service-and-container}

The service is a FastAPI application on port 8000 with
\texttt{POST\ /parse}, \texttt{GET\ /healthz} and \texttt{GET\ /readyz}.
\texttt{/parse} returns only the four-key object and sets
\texttt{X-Model-Version} and \texttt{X-Served-By} (the Pod name from a
downward-API environment variable, with \texttt{HOSTNAME} and
\texttt{socket.gethostname()} as fallbacks). The output validator checks
exact keys, the eight intents, and string-or-null slots. If the model
output is not valid, the service returns HTTP 500 with exactly
\texttt{\{"error":"invalid\_model\_output"\}}; it never returns
malformed JSON with 200. \texttt{FAULT\_INJECT=1} replaces the generated
text with \texttt{not\ json} before validation, which exercises that
path.

At start-up the service downloads
\texttt{s3://models/\textless{}version\textgreater{}/model.tar.gz} using
boto3 against \texttt{S3\_ENDPOINT}, verifies \texttt{model.sha256} if
present, safely extracts the archive, and loads the model and tokenizer.
\texttt{MODEL\_URI}, \texttt{MODEL\_VERSION}, \texttt{S3\_ENDPOINT} and
\texttt{MAX\_NEW\_TOKENS} come from the environment; in Kubernetes the
two S3 credentials come only from the \texttt{model-store-credentials}
Secret. The image is built by \texttt{service/Dockerfile} from
\texttt{python:3.12-slim}, installs the CPU-only PyTorch wheel and
pinned dependencies, and copies no weights. It runs as UID 10001 with a
read-only root filesystem and a writable \texttt{/tmp} emptyDir for the
model cache.

\texttt{\textless{}FILL:\ local\ 50-request\ result:\ p50/p95,\ peak\ RSS,\ success\ count,\ load\ time\textgreater{}}

The 50-request test was run against the local process (no Docker daemon
is available in this build environment). The container-level 1 CPU / 1
GiB test is therefore marked as not executed; the commands are
\texttt{python\ scripts/bench\_local.py\ -\/-url\ ...\ -\/-n\ 50\ -\/-pid\ \textless{}pid\textgreater{}}
and the Docker invocation in \texttt{README.md}. The service is
single-worker and generation is serialised by a lock, so peak RSS is
stable rather than accumulating.

\subsection{5. Task 4: Kubernetes deployment, scaling and
reliability}\label{task-4-kubernetes-deployment-scaling-and-reliability}

\texttt{k8s/app.yaml} creates a two-replica Deployment in \texttt{app}
with CPU request 250m/limit 1, memory request 512Mi/limit 1Gi, liveness
\texttt{/healthz} and readiness \texttt{/readyz},
\texttt{maxUnavailable:\ 0}, \texttt{maxSurge:\ 1}, a non-root read-only
filesystem, and the Secret-backed S3 credentials.
\texttt{k8s/ingress.yaml} routes \texttt{/parse}, \texttt{/healthz} and
\texttt{/readyz} through Traefik; \texttt{k8s/hpa.yaml} scales 1-2
replicas at 70\% CPU with a 60-second scale-down window.

\texttt{\textless{}FILL:\ cluster\ status.\ In\ this\ environment\ Docker/k3d\ are\ not\ installed,\ so\ the\ following\ are\ NOT\ MEASURED:\ HPA\ 1-\textgreater{}2-\textgreater{}1,\ rolling\ update\ failure\ \%,\ rollback\ header,\ pod-deletion\ recovery,\ three\ cold\ starts\ and\ median,\ per-replica\ request\ share.\ The\ exact\ commands\ are\ in\ docs/RUNBOOK.md.\textgreater{}}

Local limitation: the cluster is two k3d nodes on one laptop and shares
one CPU and one disk. Scaling to two replicas demonstrates the mechanism
but does not provide machine-level fault tolerance or a realistic
multi-node network. A real limitation of the HPA is that CPU utilisation
includes model download and start-up work; a pod can be scaled before it
is Ready, and the readiness probe is what protects the Service from
traffic. During a rolling update the old pod remains Ready until the new
pod downloads and loads the model, so the update can take as long as a
cold start.

\subsection{6. Task 5: subnetting and
NACL}\label{task-5-subnetting-and-nacl}

Requirements are allocated largest first, each immediately after the
previous block, with the smallest prefix that fits:

{\def\LTcaptype{none} % do not increment counter
\begin{longtable}[]{@{}
  >{\raggedright\arraybackslash}p{(\linewidth - 10\tabcolsep) * \real{0.1667}}
  >{\raggedright\arraybackslash}p{(\linewidth - 10\tabcolsep) * \real{0.1667}}
  >{\raggedright\arraybackslash}p{(\linewidth - 10\tabcolsep) * \real{0.1667}}
  >{\raggedright\arraybackslash}p{(\linewidth - 10\tabcolsep) * \real{0.1667}}
  >{\raggedright\arraybackslash}p{(\linewidth - 10\tabcolsep) * \real{0.1667}}
  >{\raggedright\arraybackslash}p{(\linewidth - 10\tabcolsep) * \real{0.1667}}@{}}
\toprule\noalign{}
\begin{minipage}[b]{\linewidth}\raggedright
Subnet
\end{minipage} & \begin{minipage}[b]{\linewidth}\raggedright
Network
\end{minipage} & \begin{minipage}[b]{\linewidth}\raggedright
Prefix
\end{minipage} & \begin{minipage}[b]{\linewidth}\raggedright
Usable host range
\end{minipage} & \begin{minipage}[b]{\linewidth}\raggedright
Broadcast
\end{minipage} & \begin{minipage}[b]{\linewidth}\raggedright
Usable
\end{minipage} \\
\midrule\noalign{}
\endhead
\bottomrule\noalign{}
\endlastfoot
app & 10.20.0.0 & /22 & 10.20.0.1-10.20.3.254 & 10.20.3.255 & 1022 \\
public & 10.20.4.0 & /24 & 10.20.4.1-10.20.4.254 & 10.20.4.255 & 254 \\
data & 10.20.5.0 & /25 & 10.20.5.1-10.20.5.126 & 10.20.5.127 & 126 \\
management & 10.20.5.128 & /27 & 10.20.5.129-10.20.5.158 & 10.20.5.159 &
30 \\
\end{longtable}
}

The app subnet needs 1,000 hosts, so /22 (1,022 usable) is the smallest
fit; public needs 250, so /24 (254); data needs 100, so /25 (126);
management needs 20, so /27 (30). The total consumed is 10.20.0.0/22,
10.20.4.0/24 and the two halves of 10.20.5.0/24, leaving space for
growth.

\texttt{network/nacl.json} is the public-subnet NACL. Inbound rule 100
denies 198.51.100.0/24 for all protocols and ports before the web
allows; rule 110 allows TCP 80, rule 120 allows TCP 443, rule 130 allows
TCP 22 only from 10.20.5.128/27, and rule 140 allows return traffic from
10.20.0.0/22 on ephemeral ports 1024-65535. Outbound rule 100 allows TCP
8000 to 10.20.0.0/22, and rule 110 allows replies to clients on ports
1024-65535; there is an implicit deny for everything else. The official
simulator reports \textbf{6 / 6 = 100.0\%}.

A NACL is stateless: it does not remember that an inbound connection
created a flow, so the response packets are evaluated as new packets and
need their own rules in both directions. Rule order matters because NACL
rules are evaluated from the lowest number and the first match wins; the
deny for 198.51.100.0/24 must be numbered below the TCP 80/443 allows,
otherwise those allow rules would match first and the abusive range
would be allowed.

\subsection{7. Task 6: reproducibility, testing and
CI/CD}\label{task-6-reproducibility-testing-and-cicd}

\texttt{make\ up} creates the cluster, secrets, SeaweedFS, buckets,
images, uploads the model and applies the
Deployment/Service/Ingress/HPA; \texttt{make\ down} deletes the cluster.
From a fresh clone the whole system is two commands:
\texttt{make\ install-ml} and \texttt{make\ up}, or just
\texttt{make\ up} because the target bootstraps the virtual environment
if it is missing. \texttt{make\ up} does not commit credentials:
\texttt{scripts/create\_secrets.sh} generates them, writes them to the
git-ignored \texttt{.env.local} with mode 0600, and creates the
Kubernetes Secrets.

The test suite has 49 tests covering the cleaner's rejection precedence
and duplicate rule, accounting invariant, sharding and resume, subnet
arithmetic, the official NACL simulator, service API and exact error
body, readiness before load, output validation, storage backends, and
cost arithmetic. \texttt{ruff\ check\ .} passes, \texttt{pytest\ -q}
reports \texttt{49\ passed}, and \texttt{gitleaks\ detect\ -\/-no-git}
reports \texttt{no\ leaks\ found} (evidence in \texttt{results/}).

The CI workflow in \texttt{.github/workflows/ci.yml} runs in the
required order: \texttt{lint} (ruff) -\textgreater{} \texttt{unit-tests}
(pytest) -\textgreater{} \texttt{model-evaluation} (downloads
\texttt{model.tar.gz} and \texttt{model.sha256} from the \texttt{final}
GitHub Release, verifies the checksum, runs
\texttt{training/predict.py}, then \texttt{tools/score.py})
-\textgreater{} \texttt{quality-gate} (\texttt{-\/-min-validity\ 95}). A
separate \texttt{secret-scan} job runs gitleaks. A manual
\texttt{workflow\_dispatch} input \texttt{bad\_model=true} runs the
untrained Pythia-70M through the same gate and is expected to fail; its
URL must be inserted here after the repository is pushed:

\texttt{\textless{}FILL:\ link\ to\ the\ failing\ quality-gate\ run\textgreater{}}

\texttt{AI\_USE.md} lists the AI tools and the checks applied; every
generated file was reviewed, linted, and tested, and the measured
results come from commands in \texttt{results/}, not from the assistant.

\subsection{8. Task 7: performance, cost and
review}\label{task-7-performance-cost-and-review}

\texttt{\textless{}FILL:\ local\ Locust\ results:\ p50,\ p95,\ throughput,\ error\ rate,\ total\ requests;\ note\ that\ this\ is\ a\ local\ process,\ not\ the\ cluster\textgreater{}}

The cluster measurements (\texttt{kubectl\ top\ pods}, cold-start
median, per-replica share) are not available in this build environment;
the exact collection commands are in \texttt{docs/RUNBOOK.md}. They must
be filled from a run on a Docker host before submission.

\subsubsection{8.1 Cost}\label{cost}

A replica is billed for its CPU limit and memory limit for every hour it
runs, so with two replicas at 1 vCPU and 1 GiB each:

\begin{verbatim}
per replica per hour = 1 * 3.50 + 1 * 0.45 = 3.95 INR
all replicas per hour = 2 * 3.95 = 7.90 INR
requests per hour = measured requests/second * 3600
cost per 1,000 = (7.90 / requests per hour) * 1000
\end{verbatim}

\texttt{\textless{}FILL:\ exact\ arithmetic\ with\ measured\ throughput\ and\ final\ INR\ value\textgreater{}}

\texttt{scripts/cost.py} performs exactly this calculation from
\texttt{results/loadtest\_stats.csv} and \texttt{data/rate\_card.json},
and writes \texttt{results/cost.json}.

\subsubsection{8.2 SLO}\label{slo}

SLO: \textbf{99.5\% of \texttt{/parse} requests succeed (HTTP 200) over
one load-test window}. For the measured window with \texttt{N} requests,
the error budget is \texttt{floor(0.005\ *\ N)} failed requests.
\texttt{\textless{}FILL:\ N,\ allowed\ failures,\ actual\ failures,\ budget\ used/remaining,\ met/not\ met\textgreater{}}

\subsubsection{8.3 Review}\label{review}

{\def\LTcaptype{none} % do not increment counter
\begin{longtable}[]{@{}
  >{\raggedright\arraybackslash}p{(\linewidth - 6\tabcolsep) * \real{0.2500}}
  >{\raggedright\arraybackslash}p{(\linewidth - 6\tabcolsep) * \real{0.2500}}
  >{\raggedright\arraybackslash}p{(\linewidth - 6\tabcolsep) * \real{0.2500}}
  >{\raggedright\arraybackslash}p{(\linewidth - 6\tabcolsep) * \real{0.2500}}@{}}
\toprule\noalign{}
\begin{minipage}[b]{\linewidth}\raggedright
Component or concern
\end{minipage} & \begin{minipage}[b]{\linewidth}\raggedright
Evidence (measured)
\end{minipage} & \begin{minipage}[b]{\linewidth}\raggedright
Limitation found
\end{minipage} & \begin{minipage}[b]{\linewidth}\raggedright
Improvement proposed
\end{minipage} \\
\midrule\noalign{}
\endhead
\bottomrule\noalign{}
\endlastfoot
Data and storage & \texttt{\textless{}FILL\textgreater{}} &
\texttt{\textless{}FILL\textgreater{}} &
\texttt{\textless{}FILL\textgreater{}} \\
Model quality & \texttt{\textless{}FILL\textgreater{}} &
\texttt{\textless{}FILL\textgreater{}} &
\texttt{\textless{}FILL\textgreater{}} \\
Latency and throughput & \texttt{\textless{}FILL\textgreater{}} &
\texttt{\textless{}FILL\textgreater{}} &
\texttt{\textless{}FILL\textgreater{}} \\
Resource utilisation & \texttt{\textless{}FILL\textgreater{}} &
\texttt{\textless{}FILL\textgreater{}} &
\texttt{\textless{}FILL\textgreater{}} \\
Reliability and recovery & \texttt{\textless{}FILL\textgreater{}} &
\texttt{\textless{}FILL\textgreater{}} &
\texttt{\textless{}FILL\textgreater{}} \\
Security and cost & \texttt{\textless{}FILL\textgreater{}} &
\texttt{\textless{}FILL\textgreater{}} &
\texttt{\textless{}FILL\textgreater{}} \\
\end{longtable}
}

\subsection{9. Limitations and
improvements}\label{limitations-and-improvements}

The largest limitation of this submission is environmental: the machine
used to build it has no Docker daemon and no k3d/kubectl, and its free
memory is about 1.5 GiB while \texttt{phi3.5} needs about 2.6 GiB. The
code, manifests and tests are complete, but the cluster experiments and
the full 3,000-command teacher run could not be executed here. The
report marks those values as not measured rather than inventing them.
The next step on a Docker host is to run \texttt{make\ up},
\texttt{make\ teacher-label}, \texttt{make\ clean},
\texttt{make\ train}, \texttt{make\ train-v2}, \texttt{make\ evaluate},
then the Task 4 runbook and \texttt{make\ loadtest}; each command writes
its evidence to \texttt{results/}.

Beyond the environment, a 70M-parameter student will not match a 3.8B
teacher; the smoke training data is much smaller than the required 1,500
records, so the local scores are only a pipeline check. The HPA reacts
to CPU rather than to queue depth, the service serialises generation
under a lock, and the k3d cluster is two nodes on one machine. Better
choices would be a queue-depth or latency-based autoscaler, a
batched/streaming inference server, and a multi-machine cluster to test
real node failure. On the data side, more teacher examples, stricter
slot grounding, and a held-out split would reduce overfitting. On the
security side, the S3 credentials are generated and Secret-backed, the
image is non-root and read-only, and the NACL is default deny; a network
policy restricting the app namespace to SeaweedFS and the Ingress would
add defence in depth.

\end{document}
